\documentclass[11pt,twocolumn]{article}
\usepackage[a4paper,margin=0.75in]{geometry}
\usepackage{amsmath,amssymb}
\usepackage{graphicx}
\usepackage{booktabs}
\usepackage{hyperref}
\usepackage{enumitem}
\setlength{\columnsep}{0.25in}
\documentclass[12pt,a4paper]{article}
\usepackage{placeins} % in preamble
\usepackage[section]{placeins}
% ---- packages ----
\usepackage[utf8]{inputenc}
\usepackage[T1]{fontenc}
\usepackage[margin=2.5cm]{geometry}
\usepackage{listings}
\usepackage{xcolor}
\usepackage{graphicx}
\usepackage{amsmath}
\usepackage{hyperref}
\usepackage{parskip}
\usepackage{booktabs}
\usepackage{float}
\usepackage[section]{placeins}


\title{Visualization of Set Valued Dynamical Systems}
\author{Nam Do\\University of Oulu, Finland\\\texttt{Nam.Do@student.oulu.fi} \and Kalle Timperi\\University of Oulu, Finland\\\texttt{kalle.timperi@oulu.fi} \and Jyri Ollanketo\\University of Oulu, Finland\\\texttt{jyri.ollanketo@student.oulu.fi}}
\date{April 14, 2026}

\begin{document}
\maketitle

\begin{abstract}
Set valued dynamical systems provide a rigorous framework for uncertainty propagation in nonlinear maps and flows with bounded noise. This report extends the Stage 2 implementation with continuous time dynamics, boundary differential equation tracking, a continuous Ulam Markov approximation, parameter sweeping for periodic orbits, a generic two equation parser for user defined systems, arbitrary axis range control, and a redesigned visualization interface. The software now integrates boundary tracking, unstable manifold approximation, periodic orbit detection, and Ulam transition matrix methods for both discrete maps and continuous flows. We also document the updated architecture, test suite, and continuous integration workflow, and present an evaluation focused on one key correctness item per feature based on the executed test runs.
\end{abstract}

\section{Introduction}
This report describes the design and implementation of a visualization tool for set valued dynamical systems under bounded noise. We establish the mathematical foundations for boundary tracking, unstable manifolds, periodic orbit detection, and Ulam based set oriented methods. We then present the computational implementation and the updates introduced in Stage 3, with an evaluation of correctness and reliability.

\subsection{Set Valued Dynamical Systems and Motivation}
For a deterministic map $f: \mathbb{R}^d \to \mathbb{R}^d$, bounded additive noise yields a set valued map
\begin{equation}
F(X) = \bigcup_{x \in X} B_{\epsilon}(f(x))
= \{f(x) + \xi : x \in X, \|\xi\| \le \epsilon\},
\end{equation}
where $B_{\epsilon}(y)$ is the closed $\epsilon$ ball. Set valued maps encode the worst case evolution of uncertainty and are therefore appropriate for chaotic systems where small perturbations amplify rapidly.

\subsection{Minimal Invariant Sets and Attractors}
A compact set $M \subset \mathbb{R}^d$ is invariant if $F(M) = M$, and minimal if no strict subset is invariant. For any $x \in M$,
\begin{equation}
M = \bigcup_{k=0}^{\infty} F^k(\{x\}).
\end{equation}
$M$ is an attractor if there exists $\sigma > 0$ such that
\begin{equation}
\lim_{k \to \infty} \operatorname{dist}_H(F^k(B_{\sigma}(M)), M) = 0,
\end{equation}
where $\operatorname{dist}_H$ is the Hausdorff distance. The minimal invariant set provides the smallest forward invariant region that contains all noisy trajectories.

\subsection{Dual Map and Dual Repeller}
The dual map for a diffeomorphism $f$ is the inverse $f^{-1}$. For set valued dynamics, the dual map tracks backwards reachable sets. The dual repeller is the minimal invariant set of the inverse map and provides information about topological bifurcations where attractors split or merge. We compute dual repellers using backward iteration of the boundary map and evaluate their intersection with forward invariant sets.

\subsection{Boundary Tracking Over Full Noise Balls}
Tracking the full image $F(X)$ is expensive. The maximal uncertainty occurs on the boundary of $B_{\epsilon}(f(x))$, so we evolve only boundary points and their normals. This yields the correct boundary of the minimal invariant set while avoiding interior sampling.

\subsection{Boundary Map for Discrete Systems}
Let $z \in \mathbb{R}^2$ be a boundary point and $n$ a unit outward normal. The extended boundary map is
\begin{equation}
E(z,n) = \bigl(f(z) + \epsilon n', n'\bigr), \quad
n' = \frac{(Df(z)^{-T}) n}{\|(Df(z)^{-T}) n\|}.
\end{equation}
Iterating $E$ over an initial boundary circle yields a discrete approximation to the minimal invariant boundary, following the boundary mapping construction in Tey (Section 2.2) \cite{tey2023}.


\subsection{Boundary Differential Equation for Continuous Systems}
For an ODE $\dot{x} = f(x)$, boundary evolution follows the boundary differential equation (BDE)
\begin{align}
\dot{x} &= f(x) + \epsilon n, \\
\dot{n} &= -Df(x)^T n + \bigl(Df(x)^T n \cdot n\bigr) n,
\end{align}
which preserves $\|n\| = 1$. This formulation follows Tey (Section 2.4.1) \cite{tey2023}. We integrate the BDE using a fourth order Runge Kutta scheme, and periodically reparameterize the boundary to main-
tain near uniform arc length spacing.


\subsection{Example Systems}
\paragraph{H\'enon map.} The H\'enon map is \cite{henon1976}
\begin{equation}
(x_{k+1}, y_{k+1}) = (1 - a x_k^2 + y_k, b x_k).
\end{equation}
\paragraph{Duffing map.} The Duffing map is a discrete approximation of the Duffing oscillator \cite{duffing1918},
\begin{equation}
(x_{k+1}, y_{k+1}) = (y_k, -b x_k + a y_k - y_k^3).
\end{equation}
\paragraph{Duffing ODE.} The continuous Duffing oscillator is \cite{duffing1918}
\begin{equation}
\dot{x} = y, \quad \dot{y} = x - x^3 - \delta y.
\end{equation}

\subsection{Periodic Orbit Detection}
Periodic boundary points are fixed points of $E^p$. We solve
\begin{equation}
E^p(z,n) - (z,n) = 0
\end{equation}
by Newton iteration on a grid of initial guesses, following the approach of Davidchack and Lai \cite{davidchack1999}. Stability is classified by the eigenvalues of the Jacobian $J_p = D E^p(z,n)$: stable if $|\lambda_i| < 1$, unstable if $|\lambda_i| > 1$, and saddle otherwise.

\subsection{Unstable Manifold Approximation}
We compute unstable manifolds by iterating local eigenvectors from saddle points and applying an adaptive subdivision strategy based on the subdivision method of Dellnitz and Hohmann \cite{dellnitz1997}. Given consecutive points $(z_0, n_0)$ and $(z_1, n_1)$, we interpolate a midpoint, evaluate $E$, and subdivide until the segment length is below a tolerance. This produces a refined polyline that approximates the unstable manifold and supports dual repeller visualization.

\subsection{GAIO Based Ulam Method}
We partition the phase space into a grid of boxes $\{B_i\}$. For each box $B_i$, we sample points $x_k$ and count intersections of $B_{\epsilon}(f(x_k))$ with boxes $B_j$. The transition probabilities are
\begin{equation}
P_{ij} = \frac{1}{N_i} \#\{x_k \in B_i : B_{\epsilon}(f(x_k)) \cap B_j \ne \emptyset\}.
\end{equation}
The right eigenvector of $P$ yields an invariant measure; the left eigenvector highlights dual repellers and supports bifurcation detection \cite{dellnitz2001,ulam1960}. For continuous systems, $f$ is replaced by the time $T$ flow map computed by RK4 substeps.

\subsection{Topological Bifurcation}
Bifurcations in minimal invariant sets can be detected by changes in the number of connected components and by the emergence of dual repellers \cite{lamb2015,lamb2025}. The parameter sweep interface samples a parameter and tracks periodic orbit counts and stability changes, providing an operational detection method.

\section{Implementation}
\subsection{System Architecture}
Core numerical methods are implemented in Rust and compiled to WebAssembly with \texttt{wasm-bindgen}. The frontend is a React application with Three.js based rendering. The Rust modules are organized by responsibility, including \texttt{boundary\_periodic}, \texttt{unstable\_manifold}, \texttt{ulam}, \texttt{continuous\_ds}, \texttt{user\_defined}, \texttt{parameters}, and \texttt{range}. The frontend orchestrates computation, visualization, and user interaction via the main component \texttt{SetValuedViz.jsx}.


\begin{figure}[!htbp]
  \centering
  \fbox{\includegraphics[width=\columnwidth]{stage3/system_overview.png}}
  \caption{System overview with the sidebar controls and the main dynamical system visualization viewport}
  \label{fig:task2}
\end{figure}


\subsection{Dynamical Systems Support}
The tool supports discrete H\'enon and Duffing maps, user defined discrete systems, the continuous Duffing ODE, and user defined ODEs. For user defined systems, equation strings are parsed with a safe expression evaluator and paired with a validated parameter list. Reserved identifiers such as $x$, $y$, and built in functions are rejected to avoid ambiguity. This refactor enables generic two equation systems in both discrete and continuous settings.

\subsection{Continuous Boundary Simulation}
The continuous module defines a shared \texttt{OdeSystem} trait with vector field, Jacobian, and RK4 stepping. The BDE uses a coupled RK4 step for $(x,n)$, and a reparameterization routine redistributes points by arc length to improve discretization quality. An Euler map $F_h(x) = x + h f(x)$ is included to connect continuous fixed points with discrete stability analysis.
\begin{figure}[!h]
  \centering
  \fbox{\includegraphics[width=\columnwidth]{stage3/boundary_differential_equation_visualization.png}}
  \caption{Example of Duffing equation for visualizing continuous dynamical systems}
  \label{fig:task2}
\end{figure}

\subsection{Periodic Orbits and Parameter Sweep}
Periodic orbits for the boundary map are computed with Newton iteration and stability classification. A parameter sweep interface evaluates a parameter range, aggregates orbit counts by stability type, and exports CSV or JSON. The sweep implementation supports both H\'enon and arbitrary user defined systems, and uses the same boundary periodic solver for consistency.

\subsection{Unstable and Stable Manifolds}
Unstable manifold computation follows the subdivision algorithm of Dellnitz and Hohmann, with termination criteria based on point spacing and curvature. Stable manifolds and dual repellers are computed through backward iteration of the boundary map. These elements form the skeletal structure of minimal invariant sets and are rendered as polylines alongside periodic orbits.

\subsection{Ulam and Markov Integration}
The GAIO style Ulam method is implemented for discrete maps and for continuous flows via the time $T$ map of the Duffing ODE. The grid domain is clamped to a configurable axis range, and the transition matrix is computed using epsilon inflated images. The invariant measure and dual repeller overlays are visualized as color coded grids.

\subsection{Visualization Interface and Axis Range Control}
The interface introduces a system type selector (discrete versus continuous), a parameter editor for custom equations, an axis range panel with clamping, and a parameter sweep panel. The coordinate grid adapts to user specified ranges, and the view range is normalized to avoid degenerate or inverted axes. Visualization controls allow users to toggle trajectories, periodic orbits, manifolds, and Ulam overlays.

\subsection{Documentation, Tests, and CI}
Mathematical and architectural documentation is maintained in \texttt{set-valued-viz/docs} along with implementation notes for continuous boundary tracking and parameterized systems. Automated Rust tests cover boundary maps, manifolds, continuous dynamics, parameter parsing, and range handling. Frontend tests verify UI panels and sweep logic. GitHub Actions runs \texttt{cargo build} and \texttt{cargo test} on pushes and pull requests. Presentation slides were prepared to summarize the methodology and results.

\section{Evaluation}
Each subsection reports one key evaluation item for the corresponding feature.

\subsection{Literature Consistency}
For the H\'enon map with $a = 0.4$, $b = 0.3$, and $\epsilon = 0.0625$, the boundary tracking converges to a single connected attractor consistent with the reference results in Tey's dissertation \cite{tey2022}. This case remains the primary external validation baseline for the discrete boundary map.

\subsection{Continuous BDE Consistency}
The RK4 boundary differential equation step reduces to the standard ODE RK4 step when $\epsilon = 0$. This property is verified by the unit test \texttt{test\_rk4\_bde\_step\_matches\_ode\_when\_epsilon\_zero}, which checks equality to within $10^{-12}$ in floating point norm.

\subsection{Parameter Sweep Reliability}
Generic sweep tests confirm that orbit counts are consistent across parameter samples and that CSV and JSON exports are generated with stable field ordering. The tests \texttt{test\_generic\_sweep\_orbit\_counts\_consistent} and \texttt{test\_generic\_sweep\_csv\_export} passed during the Stage 3 verification run.

\subsection{Custom Equation Validation}
User defined equations reproduce known maps. The Lozi map test evaluates $x' = 1 - a |x| + y$ and $y' = b x$ and confirms correct handling of absolute value parsing and parameter substitution.

\subsection{Axis Range Normalization}
Axis range normalization ensures bounded and non degenerate visualization domains. The Rust test \texttt{test\_clamp\_pair\_expands\_zero\_width} and the frontend test \texttt{viewRange.test.js} confirm that inverted inputs are reordered and that zero width ranges are expanded within the configured limit.

\subsection{Automated Test Execution}
All Rust and frontend tests were executed on April 14, 2026. The Rust suite reported 92 passing tests. The frontend suite reported 62 passing tests; the run emitted expected jsdom warnings for canvas context creation, but no test failures were observed.

\section{Discussion}
\subsection{Achievements}
Stage 3 extends the tool from discrete maps to continuous flows while preserving the set valued framework. The system now supports boundary differential equation tracking, continuous Ulam Markov approximations, parameter sweeping of periodic orbits, and a generic two equation parser for user defined systems. The visualization interface has been expanded with axis range controls and improved panels, and the codebase is supported by unit tests and continuous integration.

\subsection{Limitations and Future Work}
Continuous Ulam computation remains computationally expensive for large grids, and additional performance benchmarks for continuous flows are planned. The frontend test environment uses jsdom and lacks a full canvas implementation, which limits chart rendering coverage. Future work includes GPU acceleration for boundary evolution, adaptive grid refinement for Ulam matrices, and quantitative bifurcation diagrams for continuous systems.

\section{Conclusion}
We presented an extended visualization framework for set valued dynamical systems that integrates discrete and continuous dynamics with boundary tracking, manifold computation, and probabilistic Ulam methods. The new features enable exploration of ODE driven set valued dynamics, user defined models, and parameter sweeps with stable numerical foundations and improved usability. The evaluation confirms key correctness properties via automated tests and literature alignment.

\begin{thebibliography}{11}
\bibitem{tey2022} W. H. Tey, ``Boundary mappings of minimal invariant sets,'' Ph.D. dissertation, Imperial College London, 2022.
\bibitem{henon1976} M. H\'enon, ``A two dimensional mapping with a strange attractor,'' Communications in Mathematical Physics, vol. 50, no. 1, pp. 69 to 77, 1976.
\bibitem{davidchack1999} R. L. Davidchack and Y. C. Lai, ``Efficient algorithm for detecting unstable periodic orbits in chaotic systems,'' Physical Review E, vol. 60, no. 5, pp. 6172 to 6175, 1999.
\bibitem{dellnitz1997} M. Dellnitz and A. Hohmann, ``A subdivision algorithm for the computation of unstable manifolds and global attractors,'' Numerische Mathematik, vol. 75, no. 3, pp. 293 to 317, 1997.
\bibitem{dellnitz2001} M. Dellnitz, G. Froyland, and O. Junge, ``The algorithms behind GAIO, set oriented numerical methods for dynamical systems,'' in \textit{Ergodic Theory, Analysis, and Efficient Simulation of Dynamical Systems}, B. Fiedler, Ed. Springer, 2001, pp. 145 to 174.
\bibitem{ulam1960} S. M. Ulam, \textit{A Collection of Mathematical Problems}. New York: Interscience Publishers, 1960.
\bibitem{peschke2015} G. Peschke, ``Computational analysis of invariant sets and their bifurcation in set valued dynamical systems,'' M.Sc. thesis, Technische Universit\"at M\"unchen, 2015.
\bibitem{lamb2015} J. S. W. Lamb, M. Rasmussen, and C. S. Rodrigues, ``Topological bifurcations of minimal invariant sets for set valued dynamical systems,'' Proceedings of the American Mathematical Society, vol. 143, no. 9, pp. 3927 to 3937, 2015.
\bibitem{lamb2025} J. S. W. Lamb, M. Rasmussen, and W. H. Tey, ``Bifurcations of the H\'enon map with additive bounded noise,'' arXiv preprint arXiv:2504.02776, 2025.
\bibitem{duffing1918} G. Duffing, \textit{Erzwungene Schwingungen bei Ver\"anderlicher Eigenfrequenz und ihre Technische Bedeutung}. Braunschweig: Vieweg, 1918.
\bibitem{tey2023} W. H. Tey, ``Boundary mappings of minimal invariant sets,'' Ph.D. dissertation, Imperial College London, 2023.

\end{thebibliography}

\end{document}
